%% LaTeX-Beamer template for KIT design
%% by Erik Burger, Christian Hammer
%% title picture by Klaus Krogmann
%%
%% version 2.1
%%
%% mostly compatible to KIT corporate design v2.0
%% http://intranet.kit.edu/gestaltungsrichtlinien.php
%%
%% Problems, bugs and comments to
%% burger@kit.edu

\documentclass[18pt]{beamer}
\def\tempfolder{../template/}
%% SLIDE FORMAT

% use 'beamerthemekit' for standard 4:3 ratio
% for widescreen slides (16:9), use 'beamerthemekitwide'
\usepackage{\tempfolder beamerthemekit}
% \usepackage{templates/beamerthemekitwide}

% -- Settings --------------------------------------------------------%
\setbeamertemplate{navigation symbols}{}

\def\figfolder{../figs/}
\def\logofolder{\tempfolder logo/}

\newcommand{\playIcon}{%
	\begin{tikzpicture}[scale=0.5]
	\fill[fill=black] (0,0) -- (1.5,1) -- (0, 2) -- (0, 0);
	\end{tikzpicture}}


\usepackage{animate}
\usepackage{pdflscape}
\usepackage{pdfpages}
\usepackage{framed, color}
\usepackage{color}
\usepackage{tikz}
\usepackage{graphicx}
%\usepackage{subcaption}
\usepackage{listings}
\usepackage{tabulary}
\usepackage{tabularx}
\usepackage{qrcode}
\usepackage[utf8]{inputenc}

\usetikzlibrary{shapes}
\usetikzlibrary{calc}

%% TITLE PICTURE
%\titleimage{mypicture}

%% TITLE LOGO
%\titlelogo{mylogo}

%% TikZ INTEGRATION
% use these packages for PCM symbols and UML classes
% \usepackage{templates/tikzkit}
% \usepackage{templates/tikzuml}

% the presentation starts here

\title[Gruppe]{4D visualisation of the tropopause, identification of air mass exchanges and their fate}
\subtitle{Department of Mathematics \& Steinbuch Centre for Computing}
\author{Tianbai Xiao}

\institute{MathSEE Modeling Week}

%---------------------------------------------------------------------------------
\begin{document}
	
	\selectlanguage{english}
	
	%title page
	\titleimage{\tempfolder logo/background}
	\titlelogo{\tempfolder logo/MathSEE}
	%----------------------------------------------------------%
	
	%----------------------------------------------------------%
	% titlepage
	\maketitle

	%----------------------------------------------------------%
	%----------------------------------------------------------%
	\begin{frame}[c]
		\frametitle{Optimales Warenlager}
		\vspace*{-8mm}
		\begin{center}
			\begin{tikzpicture}
			\node at (0, 4)		{\includegraphics[width=0.3\textwidth]{../figs/z_laptop}};
			\node at (0, 2.7)	{\Large Kunde };
			
			\uncover<2->{
				\node at (3.6, 2)		{\includegraphics[width=0.3\textwidth]{../figs/z_lieferant}};
				\node[rotate=-35] at (1.8, 3)		{\includegraphics[width=1cm]{../figs/z_leftRightArrow}};
				\node at (3.6, 0.5)	{\Large H\"andler };
				\node at (3.6, 0)	{ (Online Shop) };
			}
			\uncover<3->{
				\node at (7.5, -0.5)		{\includegraphics[width=0.27\textwidth]{../figs/z_lager}};
				\node[rotate=-35] at (5.5, 0.5)		{\includegraphics[width=1cm]{../figs/z_leftRightArrow}};
				\node at (7.5, -2)	{\Large Lager };
			}
			\uncover<4->{
				\node [draw, rectangle, thick, fill=red, fill opacity=0.1, color=red, minimum height=5.8cm, minimum width=6.9cm] at (5.7, 0.5) {};
			}
			\end{tikzpicture}
		\end{center}
	\end{frame}
	
	
	\begin{frame}[c]
		\frametitle{Optimales Warenlager}
		\vspace*{-7mm}
		\begin{center}
			\begin{tikzpicture}
			\uncover<1->{
				\node at (3.8, -1)	{\includegraphics[width=0.25\textwidth]{../figs/z_lieferant}};
				\node at (1.8, -1.4)	{ \large H\"andler };
			}
			
			\uncover<2->{
				\node at (6.5, 2.6)	{\includegraphics[width=6cm]{../figs/z_lager}};
				\node at (6.5, 0)	{Warenlager};
			}
			
			\uncover<3->{
				\node at (6.5, 5.3)	{Artikel};
			}
			
			\uncover<3->{
				\node at (6.5, 5.3)	{Artikel};	
				\pgfmathsetmacro{\x}{3.8}
				\pgfmathsetmacro{\y}{3.7}
				\draw[draw, fill=red, opacity=0.3] ($(0,0)+(\x,\y)$) rectangle ($(0.5,0.5)+(\x,\y)$);
				
				\pgfmathsetmacro{\x}{5.7}
				\pgfmathsetmacro{\y}{3.1}
				\draw[draw, fill=darkgreen, opacity=0.6] ($(0,0)+(\x,\y)$) rectangle ($(0.4,0.4)+(\x,\y)$);
				
				\pgfmathsetmacro{\x}{8.35}
				\pgfmathsetmacro{\y}{3.2}
				\draw[draw, fill=blue, opacity=0.4] ($(0,0)+(\x,\y)$) rectangle ($(0.4,0.4)+(\x,\y)$);
				
				\pgfmathsetmacro{\x}{8.2}
				\pgfmathsetmacro{\y}{1.7}
				\draw[draw, fill=black, opacity=0.5] ($(0,0)+(\x,\y)$) rectangle ($(0.4,0.4)+(\x,\y)$);
			}
			
			\uncover<4->{
				\node at (0, 0.15)	{Kundenbestellungen};
			}
			
			\pgfmathsetmacro{\x}{0}
			\pgfmathsetmacro{\y}{0}
			\uncover<4->{
				\node[yscale=1] at ($(0, 3.5)+(\x,\y)$)	{\includegraphics[width=0.3\textwidth]{../figs/z_einkaufszettel}};
				\node at ($(0, 4.1)+(\x,\y)$)	{\large \it \tt 1 x iPhone5 };
				\node at ($(0, 3.4)+(\x,\y)$)	{\large \it \tt 3 x Rasier- };
				\node at ($(0.5, 3.0)+(\x,\y)$)	{\large \it \tt klingen };
				\node at ($(0, 2.4)+(\x,\y)$)	{\large \it \tt 5 x Stifte };
				%\draw[draw, fill=red, opacity=0.2] ($(-1.65,1.3)+(\x,\y)$) rectangle ($(1.65,5.7)+(\x,\y)$);
			}
			
			\pgfmathsetmacro{\x}{0.35}
			\pgfmathsetmacro{\y}{-0.35}
			\uncover<5->{
				\node[yscale=1] at ($(0, 3.5)+(\x,\y)$)	{\includegraphics[width=0.3\textwidth]{../figs/z_einkaufszettel}};
				\node at ($(0, 4.1)+(\x,\y)$)	{\large \it \tt 3 x Socken };
				\node at ($(0, 3.4)+(\x,\y)$)	{\large \it \tt 2 x Heftordner };
				\node at ($(0.5, 3.0)+(\x,\y)$)	{\large \it \tt  };
				\draw[draw, fill=red, opacity=0.15] ($(-1.65,1.3)+(\x,\y)$) rectangle ($(1.65,5.7)+(\x,\y)$);
			}
			
			\pgfmathsetmacro{\x}{0.7}
			\pgfmathsetmacro{\y}{-0.7}
			\uncover<6->{
				\node[yscale=1] at ($(0, 3.5)+(\x,\y)$)	{\includegraphics[width=0.3\textwidth]{../figs/z_einkaufszettel}};
				\node at ($(0, 4.1)+(\x,\y)$)	{\large \it \tt 1 x Dschungel-};
				\node at ($(0.5, 3.75)+(\x,\y)$)	{\large \it \tt buch DVD };
				\node at ($(0, 3.1)+(\x,\y)$)	{\large \it \tt 1 x Teddybär};
				\node at ($(0, 2.5)+(\x,\y)$)	{\large \it \tt 2 x Handschuhe};
				\draw[draw, fill=darkgreen, opacity=0.4] ($(-1.65,1.3)+(\x,\y)$) rectangle ($(1.65,5.7)+(\x,\y)$);
			}
			
			\end{tikzpicture}
		\end{center}
	\end{frame}
	
	
	\begin{frame}[c]
		\frametitle{Optimales Warenlager}
		\large
		
		{\bf Gegeben}
		\begin{itemize}
			\item Lageranordnung (Gänge, Regale, etc.)
			\item Lagerartikel
			\item Kundenbestellungen
		\end{itemize}
		
		\vspace{5mm}
		
		\uncover<2->{
			{\bf Gesucht}
			\begin{itemize}
				\item Anordnung der Artikel im Lager?
				\item Ziel: Minimaler Weg
			\end{itemize}
		}
	\end{frame}
	
\end{document}
% ---------------------------------------------------------------------------------
